\section{Conclusão}

\subsection{Síntese dos achados}

Os quatro objetivos do trabalho foram endereçados com os dados reais do efluente da LAGWRP (2011--2026, 182 meses), com uma bateria de 21 métodos de previsão (10 originais + 4 \textit{baselines} obrigatórios + 7 métodos adicionais), todos validados com MASE, validação cruzada expansiva e \textit{backtesting} de origem móvel, além do \textit{holdout} tradicional. (1) A tendência de TDS sobre a série completa é de alta e estatisticamente significativa: cinco métodos com ajuste global à série inteira (Mann-Kendall/Sen, Theil-Sen, OLS, STL+tendência e regressão bayesiana) convergem para uma taxa de aproximadamente 2,9 a 3,9 mg/L/ano (p $<$ 0,01 em todos os casos). Esse resultado, porém, não é robusto a duas escolhas metodológicas centrais: restrita ao período posterior à troca de código do ponto de monitoramento (\texttt{EFF-001A}, 2012--2026, 170 dos 182 meses), a inclinação cai para 0,84 mg/L/ano e deixa de ser significativa (p=0,59); e a significância do teste de Mann-Kendall se inverte conforme a correção de autocorrelação serial adotada (significativa sob a correção padrão, sob \textit{trend-free pre-whitening} e sob Seasonal Kendall; não significativa sob \textit{pre-whitening} simples e sob a correção de variância de Hamed e Rao). Investigação dedicada mostrou que a transição de código de monitoramento em si não introduz um degrau artificial (transição contínua mês a mês, mesmo método analítico, mesmo limite de detecção, mesmas coordenadas) -- em vez disso, a série de TDS se comporta de forma cíclica por regime, não como uma reta ascendente única: alta de 563 para 808 mg/L entre 2011 e 2015 (coincidindo com a seca da Califórnia de 2012--2016), queda até 636 mg/L em 2019, nova alta até 768 mg/L em 2022 (segunda seca, 2020--2022) e leve recuo depois. A vazão do efluente, reconstruída a partir da identidade $\text{lb/dia} = \text{mg/L} \times \text{vazao(MGD)} \times 8,34$ e validada por consistência cruzada entre quatro parâmetros medidos na mesma amostra física, correlaciona-se com o TDS no sentido esperado pela hipótese de diluição (r anual $=-0{,}55$, p=0,027), mas não reproduz integralmente o formato fino desse ciclo -- entre 2016 e 2017, por exemplo, a vazão permaneceu estável enquanto o TDS caiu de forma acentuada. A leitura mais honesta dos dados deste trabalho é, portanto, que a série de TDS da LAGWRP responde a um padrão cíclico plausivelmente ligado a ciclos de seca, e não a uma tendência monotônica de longo prazo -- ver Seção "Tratamento de dados robusto" dos Resultados para a matriz completa de sensibilidade. (2) A previsão para +10/+15/+20 anos não tem um valor único confiável, mas os três finalistas recomendados (regressão bayesiana, Detrend+RF e híbrido SARIMA+Prophet -- Seção "Métodos adicionais"/"Diagnóstico de resíduos") convergem para um intervalo razoavelmente estreito de ~760 a ~800 mg/L em +20 anos, apesar de mecanismos de extrapolação bem diferentes entre si. Essa convergência é evidência de curto/médio prazo, mas a Seção "WRTDS" e a Seção "Projeção por cenários climáticos" dos Resultados mostram que ela carrega implicitamente a suposição de que a vazão continua caindo -- um pressuposto que não se sustenta fisicamente além de ~10 anos (Seção "Balanço de massa"). A limitação de extrapolação dos métodos baseados em árvore originais (Random Forest, XGBoost, LightGBM, que saturam ou produzem "tendências" implícitas negativas) foi corrigida pelo método Detrend+árvore (Seção "Métodos adicionais"), que separa a tendência (extrapolada linearmente) do resíduo (modelado pela árvore, sem precisar extrapolar); SARIMA e Prophet continuam divergindo em direções opostas ao extrapolar isoladamente, mas o ensemble híbrido SARIMA+Prophet atenua parcialmente essa divergência. O diagnóstico de resíduos (Ljung-Box, Shapiro-Wilk, ARCH) mostrou que o Detrend+RF é o único candidato sem autocorrelação residual nem heterocedasticidade condicional detectável, reforçando sua indicação como finalista. (3) A correlação entre TDS e Amônia é positiva e estatisticamente significativa mesmo após remover a tendência comum ($r=0{,}175$, p=0,018; mais forte ainda com defasagem de 1 mês, $r=0{,}215$), oferecendo evidência de que meses com salinidade acima do esperado tendem a coincidir com desempenho de nitrificação abaixo do esperado. Já a correlação entre TDS e BOD é \textbf{negativa} e não estatisticamente significativa nos três tratamentos de ND testados ($r$ entre $-0{,}068$ e $-0{,}126$; p entre 0,090 e 0,360, incluindo o ROS/Helsel, o mais bem embasado estatisticamente) -- sinal oposto ao previsto pela hipótese de que TDS alto prejudica a remoção de matéria orgânica. Esse resultado não confirma a hipótese do professor e é reportado como tal, sem suavização; a leitura mais provável, dada a robustez do sinal entre tratamentos, é a falta de variação real na série de BOD (que permanece rente ao limite de detecção do método durante quase todo o período) mais do que um efeito genuíno na direção oposta. (4) Cloreto como regressor exógeno adicional (SARIMAX) tem coeficiente estatisticamente significativo mas não melhora de forma prática a previsão de TDS sobre o SARIMA univariado -- resultado negativo também reportado, não omitido. (5) A confirmação do padrão cíclico (D-14/D-37) motivou três métodos desenhados especificamente para essa estrutura, respondendo diretamente à pergunta causal do professor: o \textbf{WRTDS} mostra que a tendência flow-normalized (o que resta depois de descontar o efeito da vazão) não é estatisticamente significativa ($p=0{,}064$, robusto à checagem de circularidade) -- quase toda a tendência bruta é efeito de vazão, não de mais sal entrando no sistema. O \textbf{balanço de massa} confirma essa leitura diretamente: a carga de sal caiu $-3{,}20\%$/ano e a vazão caiu $-3{,}85\%$/ano no mesmo período ($p<0{,}0001$ nos dois) -- \textbf{o mecanismo é diluição, não mais sal}, mas a extrapolação linear desses dois componentes por 20 anos cruza valores fisicamente impossíveis (vazão negativa), expondo a fragilidade de qualquer previsão pontual de longo prazo nesta série. A \textbf{projeção por cenários climáticos} substitui essa previsão pontual por uma faixa condicionada ao regime de seca (594 a 721~mg/L em +20 anos, dependendo do cenário) -- mais estreita e mais baixa que a convergência dos finalistas de (2) porque limita a vazão futura a valores historicamente plausíveis, em vez de extrapolá-la livremente.

\subsection{Interpretação em contexto real}

Os achados são compatíveis com o mecanismo descrito por Schwabe et al.\ \cite{schwabe2020unintended}: medidas de conservação de água reduzem o volume de água que entra no sistema de esgoto, e, mantendo-se a mesma carga de sais, a concentração de TDS no esgoto tende a subir. O padrão cíclico de TDS identificado neste trabalho (Seção "Síntese dos achados") é consistente com essa hipótese e também com o achado de Wolfand et al.\ \cite{wolfand2022dilution}, que modelaram a mesma bacia do rio Los Angeles (onde a LAGWRP descarrega) e encontraram que a concentração de TDS aumenta com o reúso de água, mesmo quando a carga total de poluentes diminui -- mas é importante ser preciso sobre o que os dados aqui permitem afirmar, e sobre o peso relativo desse mecanismo. Um relatório técnico da Southern California Salinity Coalition \cite{scsc2018tds}, que decompõe estatisticamente a variabilidade de TDS de 26 estações de tratamento da mesma região usando o método LMG (\textit{relative importance decomposition}), atribui a maior parte da variabilidade do TDS de influente (~88\%, contra ~12\% do consumo per capita) ao TDS da própria água de abastecimento -- que varia com a origem da água importada (State Water Project, Colorado River Aqueduct) em resposta a ciclos de seca, e não está disponível como série no dataset usado neste trabalho. Isso não invalida o mecanismo de conservação de Schwabe et al., mas sugere que ele é, na melhor leitura da literatura disponível, um efeito secundário e não o determinante dominante da variabilidade observada -- apresentá-lo como \textit{o} mecanismo seria impreciso. O balanço de massa deste trabalho (Seção "Balanço de massa" dos Resultados) dá uma resposta direta e própria a essa questão: a carga de sal no efluente da LAGWRP não está subindo -- está caindo ($-3{,}20\%$/ano) --, e a vazão está caindo mais rápido ainda ($-3{,}85\%$/ano). Isso é evidência local e concreta de que o mecanismo em operação na LAGWRP é diluição (menos água diluindo uma carga de sal estável ou decrescente), não mais sal entrando no sistema -- alinhado com o mecanismo de Schwabe et al., mas obtido de forma independente, sem depender de replicar a magnitude relatada por aquele estudo. Este estudo não inclui uma série temporal de consumo de água, de adoção de medidas de conservação, nem do TDS da água de abastecimento em Los Angeles, então o padrão de TDS observado é compatível com os mecanismos descritos nesses trabalhos, não uma prova causal direta de qual deles predomina. A comparação com ambos os estudos é qualitativa (concordância na direção do efeito), não numérica -- o texto completo de nenhum dos dois foi acessível nesta sessão (Seção "Comparação com a literatura" dos Resultados), e nenhuma magnitude desses trabalhos é reproduzida aqui sem verificação direta. A correlação positiva entre TDS e amônia fortalece a plausibilidade biológica do mecanismo de preocupação levantado pelo professor (salinidade inibindo processos biológicos do tratamento): amônia residual mais alta em meses de TDS mais alto é o tipo de sinal que se esperaria observar se a nitrificação estiver, de fato, sendo prejudicada pela salinidade crescente.

Para o planejamento de infraestrutura e gestão ambiental da LAGWRP, isso sugere que monitorar a trajetória de TDS não é apenas uma questão regulatória isolada, mas potencialmente um indicador antecedente de estresse sobre a etapa biológica do tratamento (nitrificação em particular). A ausência de correlação com BOD não deve ser lida como "o BOD não é afetado" -- a série de BOD observada tem pouquíssima variação para detectar qualquer efeito, o que é, em si, mais uma limitação dos dados disponíveis do que uma conclusão sobre o processo físico, e essa conclusão se mantém mesmo com o tratamento de ND mais bem embasado estatisticamente (ROS/Helsel) disponível. Informações institucionais adicionais sobre a LAGWRP, o programa de água reciclada (OWLA) da LASAN e as políticas de conservação de água da LADWP (indicadas pelo professor como material de apoio) fornecem o contexto operacional e de política pública dentro do qual esses números devem ser lidos, ainda que não tenham sido incorporadas como dados quantitativos a este trabalho.

\subsection{Limitações}

(a) O estudo é de um único ponto de descarte de uma única estação, não generalizável sem cautela a outras ETEs; (b) a análise de correlação é observacional, não estabelece causalidade -- vale mesmo quando uma variável climática (seca) é usada para explicar os ciclos de TDS: associação temporal compatível com um mecanismo não é prova causal, e o desenho deste trabalho não inclui controle nem intervenção; (c) a série de BOD é fortemente censurada (65\% dos meses ND) e tem amplitude muito pequena mesmo quando detectada, limitando o poder estatístico da correlação TDS-BOD -- mesmo a técnica de imputação mais bem embasada estatisticamente testada (ROS/Helsel) produz um único valor substituto para as 118 observações censuradas, sem recuperar variação mês a mês real; (d) apenas 1 dos 21 métodos (SVR) atinge $R^2$ positivo no \textit{holdout} de 24 meses, evidenciando que flutuações de curto prazo não são bem capturadas pela maioria dos métodos desta bateria -- e o próprio \textit{baseline} naive supera a maioria dos métodos sofisticados nesse critério específico de curto prazo (Seção "Validação honesta e baselines obrigatórios"), embora não capture tendência de longo prazo alguma; (e) SARIMA e Prophet isolados continuam pouco confiáveis para extrapolação além do histórico de treino (o híbrido SARIMA+Prophet atenua, não elimina, essa divergência), e o Gaussian Process com \textit{kernel} composto teve o pior desempenho de toda a bateria ao reverter à média em vez de extrapolar -- limitações reportadas explicitamente, não ocultadas; (f) o LSTM leve testado não superou os métodos clássicos/\textit{baselines}, confirmando para esta série a expectativa da literatura consultada; (g) a extrapolação do Cloreto futuro (necessária para o SARIMAX multivariado nos horizontes de +10/+15/+20 anos) usa uma tendência simples própria cuja incerteza não é propagada ao IC90 final do TDS, subestimando a incerteza real desse método; (h) variáveis potencialmente relevantes (temperatura, precipitação, volume de água tratada, mudanças operacionais na estação) não foram incorporadas como covariáveis; (i) a significância estatística da tendência de TDS depende da correção de autocorrelação serial adotada no teste de Mann-Kendall -- duas das quatro correções testadas (\textit{pre-whitening} simples e correção de variância de Hamed e Rao) tornam o resultado não significativo, o que por si só já recomendaria cautela antes mesmo da reformulação em torno do padrão cíclico; (j) a variável que a literatura da região aponta como a mais explicativa da variabilidade de TDS -- o TDS da própria água de abastecimento, responsável por cerca de 88\% da variabilidade segundo o relatório técnico da Southern California Salinity Coalition \cite{scsc2018tds} -- não está disponível no dataset usado neste trabalho, de modo que os modelos aqui construídos explicam a série sem acesso ao seu principal determinante hipotético; (k) a bateria original de previsão de +10/+15/+20 anos (Seções "Bateria de metodologias" e "Diagnóstico de resíduos") foi construída e validada assumindo majoritariamente uma dinâmica de tendência (ainda que com incerteza crescente), não o padrão cíclico por regime identificado depois -- essa lacuna foi endereçada por 8 métodos adicionais desenhados especificamente para o padrão cíclico (WRTDS, balanço de massa, cenários, espaço de estados, GAM, regressão quantílica, análise de intervenção, modelo fundacional -- achado 5 da Síntese, D-39 a D-49), mas os 21 métodos originais em si não foram re-treinados sob o novo enquadramento; (l) a implementação de WRTDS é própria, não o pacote de referência \texttt{EGRET} (sem equivalente Python maduro), com janelas de meia-largura fixas em vez da expansão adaptativa do pacote real -- simplificação declarada, não uma reimplementação completa do método original de Hirsch et al.; (m) a simulação de cenários climáticos propaga a incerteza do coeficiente da regressão PDSI$\to$TDS e do próprio PDSI simulado, mas não a incerteza de reconstrução do PDSI histórico (erro de medição/interpolação do nClimDiv), o que provavelmente subestima a largura real das faixas reportadas; (n) o GAM exigiu um piso declarado na suavização do termo de tendência para evitar extrapolação explosiva -- a escolha de $\lambda\geq1{,}0$ é defensável mas não única, e um piso diferente produziria uma previsão de longo prazo diferente; (o) a regressão quantílica apresenta cruzamento de quantis nos três horizontes de previsão (Q90 previsto abaixo de Q50 em +20a), limitação estrutural de ajustar quantis de forma independente sem restrição de monotonicidade; (p) o termo sazonal MA do SARIMAX com regressores de evento mostrou instabilidade de estimação (convergiu para o limite do espaço admissível), o que recomenda cautela na leitura da incerteza dessa componente específica, embora não invalide os coeficientes dos eventos; (q) o modelo fundacional \textit{zero-shot} testado (Chronos-Bolt-Small) não superou o \textit{baseline} naive, resultado esperado e declarado antes de rodar, mas isso significa que TimeGPT e TimesFM -- não testados por indisponibilidade de API paga e por decisão de escopo, respectivamente -- permanecem como possibilidade não descartada empiricamente, apenas não priorizada.

\subsection{Trabalhos futuros}

Direções que decorrem diretamente das limitações acima e do que ainda não foi possível verificar nesta sessão: (i) modelagem multivariada incorporando volume de água tratada/consumo residencial de água como covariável direta, para testar o mecanismo de Schwabe et al.\ de forma mais próxima da causal; (ii) tratamento do BOD como dado censurado por regressão explícita (tipo Tobit) ou por Kaplan-Meier para as estatísticas-resumo (complementar ao ROS já implementado, que resolve apenas a substituição pontual necessária para a série mensal -- Antweiler e Taylor \cite{antweiler2008evaluation} recomendam Kaplan-Meier especificamente para esse caso); (iii) propagação formal da incerteza do Cloreto extrapolado para o IC90 do SARIMAX multivariado (ex.\ via simulação de Monte Carlo em vez do valor pontual usado aqui); (iv) comparação numérica com outras ETEs da região (ex.\ via o relatório da Southern California Salinity Coalition \cite{scsc2018tds}, cujas 26 estações não incluem a LAGWRP nem o Tillman -- D-22 --, ou acesso institucional ao texto completo de Schwabe et al.\ e Wolfand et al.) para avaliar se a magnitude do mecanismo de diluição aqui encontrado é típica ou atípica -- buscas por dados internacionais baixáveis de TDS de efluente foram feitas e não localizaram nenhum dataset utilizável (\texttt{material\_apoio\_referencias.md} Tema 10), então essa comparação segue necessariamente qualitativa; (v) reimplementação do WRTDS via o pacote R \texttt{EGRET} (janela adaptativa completa) para verificar se a simplificação de janela fixa usada aqui (\texttt{script\_19}) altera a conclusão de que a tendência \textit{flow-normalized} não é significativa; (vi) re-treinar a bateria de tendência/ML original (\texttt{script\_01}--\texttt{script\_15}) diretamente sob o enquadramento cíclico/de regime confirmado (D-14/D-37), em vez de assumir tendência monotônica a priori -- os 5 métodos adicionados na Etapa 3 (espaço de estados, GAM, regressão quantílica, análise de intervenção, modelo fundacional) já respondem parte dessa pergunta, mas nenhum re-treina os 21 métodos originais sob o novo enquadramento (limitação (k) acima).
